\documentclass[11pt, a4paper]{article}
\usepackage[a4paper, margin=2.5cm]{geometry}
\usepackage{fontspec}
\usepackage[spanish, bidi=basic, provide=*]{babel}
\babelprovide[import, onchar=ids fonts]{spanish}
\babelfont{rm}{Noto Sans}
\usepackage{enumitem}
\setlist[itemize]{label=-}
\usepackage{amsmath}
\usepackage{booktabs}
\usepackage{tikz}

\begin{document}

\begin{center}
    \Large \textbf{Resolución Detallada - Ejercicio 104}
\end{center}

\vspace{0.5cm}

\textbf{Enunciado:}
\begin{itemize}
    \item[a.] Hallar la distribución de probabilidad de niños y niñas en familias de 3 hijos, suponiendo iguales probabilidades de niños y niñas.
    \item[b.] Representar gráficamente la distribución en (a).
    \item[c.] Hallar la función de distribución $F(x)$ y representar gráficamente.
\end{itemize}

\vspace{0.5cm}

\section*{Parte a) Distribución de probabilidad}

Definimos nuestra variable aleatoria discreta. Sea $X$: ``cantidad de niños (varones) en una familia de 3 hijos''. 

Dado que la probabilidad de que nazca un niño o una niña es igual, tenemos que la probabilidad de éxito (niño) es $p = 0.5$. El espacio muestral para 3 nacimientos (donde V=Varón y M=Mujer) tiene $2^3 = 8$ resultados posibles, todos equiprobables:
\[ \Omega = \{MMM, VMM, MVM, MMV, VVM, VMV, MVV, VVV\} \]

Calculamos la probabilidad $f(x) = P(X=x)$ para cada valor posible de $X$ (0, 1, 2, o 3 niños):
\begin{itemize}
    \item \textbf{$X=0$ (0 niños, 3 niñas):} Solo hay 1 caso (MMM). $f(0) = \frac{1}{8}$
    \item \textbf{$X=1$ (1 niño, 2 niñas):} Hay 3 casos (VMM, MVM, MMV). $f(1) = \frac{3}{8}$
    \item \textbf{$X=2$ (2 niños, 1 niña):} Hay 3 casos (VVM, VMV, MVV). $f(2) = \frac{3}{8}$
    \item \textbf{$X=3$ (3 niños, 0 niñas):} Solo hay 1 caso (VVV). $f(3) = \frac{1}{8}$
\end{itemize}

La tabla de la función de probabilidad (o distribución de probabilidad) es:
\begin{center}
\begin{tabular}{c c c c c}
\toprule
$X$ & $0$ & $1$ & $2$ & $3$ \\
\midrule
$f(x) = P(X=x)$ & $1/8$ & $3/8$ & $3/8$ & $1/8$ \\
\bottomrule
\end{tabular}
\end{center}

\vspace{0.4cm}
\textbf{Segunda opción: Pensando los casos con Combinatoria} \\
En lugar de enumerar y contar los elementos del espacio muestral uno por uno, podemos calcular los casos posibles y favorables usando técnicas de conteo. Los \textbf{casos posibles} se calculan como variaciones con repetición (2 sexos posibles en 3 nacimientos): $VR_{2,3} = 2^3 = 8$.

Para los \textbf{casos favorables}, deseamos ubicar a $x$ niños en los 3 nacimientos. La cantidad de formas de elegir esos $x$ lugares, sin importar el orden de selección, está dada por las \textbf{combinaciones} $C_{3,x} = \binom{3}{x}$. Por lo tanto:
\begin{itemize}
    \item Para $X=0$: Elegimos $0$ lugares para niños: $\binom{3}{0} = 1$ caso favorable $\rightarrow P(X=0) = \frac{1}{8}$
    \item Para $X=1$: Elegimos $1$ lugar para niño: $\binom{3}{1} = 3$ casos favorables $\rightarrow P(X=1) = \frac{3}{8}$
    \item Para $X=2$: Elegimos $2$ lugares para niños: $\binom{3}{2} = 3$ casos favorables $\rightarrow P(X=2) = \frac{3}{8}$
    \item Para $X=3$: Elegimos $3$ lugares para niños: $\binom{3}{3} = 1$ caso favorable $\rightarrow P(X=3) = \frac{1}{8}$
\end{itemize}

\vspace{0.5cm}

\section*{Parte b) Representación gráfica de $f(x)$}

Para graficar la función de probabilidad puntual $f(x)$, simplemente levantamos una barra (o bastón) en cada valor de la variable $X$, cuya altura corresponde a su probabilidad.

\begin{center}
\begin{tikzpicture}[scale=1.5]
    % Ejes
    \draw[->, thick] (-0.5,0) -- (4,0) node[right] {$x$};
    \draw[->, thick] (0,-0.1) -- (0, 1.2) node[above] {$f(x)$};
    
    % Marcas en X
    \foreach \x in {0,1,2,3} {
        \draw (\x, -0.05) -- (\x, 0.05) node[below=0.1cm] {\x};
    }
    
    % Marcas en Y
    \draw (-0.05, 0.375) -- (0.05, 0.375) node[left=0.1cm] {1/8};
    \draw (-0.05, 1.125) -- (0.05, 1.125) node[left=0.1cm] {3/8};
    
    % Barras
    \draw[line width=12pt, blue, opacity=0.7] (0,0) -- (0,0.375);
    \draw[line width=12pt, blue, opacity=0.7] (1,0) -- (1,1.125);
    \draw[line width=12pt, blue, opacity=0.7] (2,0) -- (2,1.125);
    \draw[line width=12pt, blue, opacity=0.7] (3,0) -- (3,0.375);
\end{tikzpicture}
\end{center}

\newpage

\section*{Parte c) Función de Distribución Acumulada $F(x)$}

La función de distribución $F(x)$ se define como la probabilidad de que la variable aleatoria tome un valor menor o igual a un cierto $x$. Es decir: $F(x) = P(X \le x)$. Para calcularla, vamos ``acumulando'' (sumando) las probabilidades de la tabla anterior de izquierda a derecha.

\begin{itemize}
    \item Si $x < 0$: Todavía no hemos acumulado ninguna probabilidad. $F(x) = 0$.
    \item Si $0 \le x < 1$: Acumulamos la probabilidad de $X=0$. $F(x) = 1/8$.
    \item Si $1 \le x < 2$: Acumulamos $X=0$ y $X=1$. $F(x) = 1/8 + 3/8 = 4/8 = 1/2$.
    \item Si $2 \le x < 3$: Acumulamos $X=0, 1, y 2$. $F(x) = 1/2 + 3/8 = 7/8$.
    \item Si $x \ge 3$: Acumulamos todo el espacio muestral. $F(x) = 7/8 + 1/8 = 1$.
\end{itemize}

La definición formal (analítica) queda expresada mediante una función por partes:
\[
F(x) = 
\begin{cases} 
0 & \text{si } x < 0 \\
1/8 & \text{si } 0 \le x < 1 \\
1/2 & \text{si } 1 \le x < 2 \\
7/8 & \text{si } 2 \le x < 3 \\
1 & \text{si } x \ge 3 
\end{cases}
\]

\subsection*{Paso a Paso: Cómo interpretar y construir el gráfico de $F(x)$}

Las funciones de distribución acumulada para variables aleatorias discretas siempre tienen forma de ``escalera''. Analicemos sus elementos particulares que suelen causar confusión:

\begin{enumerate}
    \item \textbf{Tramos horizontales (los ``escalones''):} La función es plana entre número entero y número entero. Por ejemplo, la probabilidad acumulada de tener $1.5$ niños es exactamente la misma que tener $1$ niño, porque es imposible tener medio niño. Mientras no lleguemos al siguiente valor posible de $X$ (que es $2$), la suma de probabilidades no crece.
    \item \textbf{Los saltos (la altura del escalón):} Justo cuando llegamos a un valor válido de $X$ (ej. $x=1$), la función ``salta'' hacia arriba. La medida exacta de ese salto es la probabilidad puntual de ese número. En $x=1$, el gráfico salta desde $1/8$ hasta $1/2$. Ese salto mide exactamente $3/8$, que es $P(X=1)$.
    \item \textbf{Líneas punteadas verticales:} \textbf{¡Ojo aquí!} Estas líneas \textbf{no} forman parte de la función real. Se dibujan únicamente como una ayuda visual para ver dónde se conecta el final de un escalón con el inicio del siguiente, y para evidenciar el tamaño del salto. Si fueran una línea sólida continua, violarían la definición matemática de función (un valor de $x$ tendría muchos valores de $y$).
    \item \textbf{Puntos llenos y vacíos (extremos de los intervalos):}
    \begin{itemize}
        \item \textbf{Punto lleno (cerrado):} Indica el valor exacto de la función en ese punto. Se ubican en el extremo izquierdo de cada escalón (porque nuestra definición dice $x \ge k$).
        \item \textbf{Punto vacío (abierto):} Indica que la línea horizontal llega hasta ahí, pero justo en ese borde exacto, el valor ya no pertenece a ese escalón sino al de arriba (porque nuestra definición dice $x < k+1$).
    \end{itemize}
\end{enumerate}

\vspace{0.5cm}
\begin{center}
\begin{tikzpicture}[scale=1.8]
    % Ejes
    \draw[->, thick] (-1.5,0) -- (4.5,0) node[right] {$x$};
    \draw[->, thick] (0,-0.1) -- (0, 1.2) node[above] {$F(x)$};
    
    % Cuadrícula auxiliar y marcas X
    \foreach \x in {-1,0,1,2,3,4} {
        \draw (\x, -0.05) -- (\x, 0.05) node[below=0.1cm] {\x};
    }
    
    % Marcas Y (Probabilidades acumuladas con decimales)
    \foreach \y/\ytext in {0.125/0.125, 0.5/0.5, 0.875/0.875, 1/1} {
       \draw (-0.05,\y) -- (0.05,\y) node[left=0.1cm] {\footnotesize $\ytext$};
    }
    
    % Segmentos de la función escalonada (líneas continuas)
    \draw[very thick, blue] (-1.5, 0) -- (0, 0);
    \draw[very thick, blue] (0, 0.125) -- (1, 0.125);
    \draw[very thick, blue] (1, 0.5) -- (2, 0.5);
    \draw[very thick, blue] (2, 0.875) -- (3, 0.875);
    \draw[very thick, blue] (3, 1) -- (4.2, 1);
    
    % Líneas punteadas (NO son parte de la función, solo guías visuales)
    \draw[dashed, gray] (0, 0) -- (0, 0.125);
    \draw[dashed, gray] (1, 0.125) -- (1, 0.5);
    \draw[dashed, gray] (2, 0.5) -- (2, 0.875);
    \draw[dashed, gray] (3, 0.875) -- (3, 1);
    
    % Puntos vacíos (límite no incluido)
    \filldraw[fill=white, draw=blue, thick] (0,0) circle (1.5pt);
    \filldraw[fill=white, draw=blue, thick] (1,0.125) circle (1.5pt);
    \filldraw[fill=white, draw=blue, thick] (2,0.5) circle (1.5pt);
    \filldraw[fill=white, draw=blue, thick] (3,0.875) circle (1.5pt);
    
    % Puntos llenos (valor exacto en el salto)
    \filldraw[fill=blue, draw=blue, thick] (0,0.125) circle (1.5pt);
    \filldraw[fill=blue, draw=blue, thick] (1,0.5) circle (1.5pt);
    \filldraw[fill=blue, draw=blue, thick] (2,0.875) circle (1.5pt);
    \filldraw[fill=blue, draw=blue, thick] (3,1) circle (1.5pt);
    
    % Anotaciones explicativas en el gráfico
    \node[anchor=west, text=red, font=\footnotesize] at (1.1, 0.3) {Salto = $P(X=1) = 3/8$};
    \draw[->, red, shorten >=2pt] (1.1, 0.3) -- (1, 0.3125);

\end{tikzpicture}
\end{center}

\vspace{0.5cm}

\section*{Ítem Extra: Cálculo mediante la Distribución Binomial}

Aunque el enunciado original se resolvió desarrollando el espacio muestral completo, este problema se ajusta perfectamente al modelo de una \textbf{Distribución Binomial}. 

\subsection*{Supuestos de la Distribución Binomial y su cumplimiento}
Para que una variable aleatoria siga una distribución binomial, el experimento debe cumplir cuatro condiciones fundamentales (ensayos de Bernoulli):
\begin{enumerate}
    \item \textbf{Número fijo de ensayos ($n$):} En este caso, la cantidad de repeticiones del experimento base está fijada en $n = 3$ (los tres hijos de la familia).
    \item \textbf{Resultados dicotómicos (éxito/fracaso):} Cada observación (nacimiento) se clasifica en solo dos categorías mutuamente excluyentes (niño o niña). Definimos nuestro ``éxito'' estadístico como el nacimiento de un niño, ya que la variable $X$ cuenta la cantidad de niños.
    \item \textbf{Probabilidad de éxito constante ($p$):} El enunciado asume iguales probabilidades para niño y niña, por lo que la probabilidad de que nazca un niño es $p = 0.5$. Esta probabilidad permanece constante para el primer, segundo y tercer hijo.
    \item \textbf{Independencia:} El resultado (sexo) de un nacimiento no afecta de ninguna manera la probabilidad del sexo de los nacimientos posteriores.
\end{enumerate}

Como se cumplen todos los supuestos de manera rigurosa, podemos afirmar que la variable aleatoria $X$ se distribuye de forma binomial con parámetros $n=3$ y $p=0.5$, lo cual se denota matemáticamente como $X \sim B(3; 0.5)$.

\subsection*{Cálculo de las probabilidades}
La fórmula de la función de probabilidad para la distribución binomial nos permite evitar el desarrollo de todo el espacio muestral:
\[ f(x) = P(X=x) = \binom{n}{x} p^x (1-p)^{n-x} \]

Aplicando nuestros parámetros, calculamos las probabilidades para cada valor del recorrido $R_X = \{0, 1, 2, 3\}$:
\begin{itemize}
    \item $P(X=0) = \binom{3}{0} (0.5)^0 (1-0.5)^{3} = 1 \cdot 1 \cdot 0.125 = \frac{1}{8}$
    \item $P(X=1) = \binom{3}{1} (0.5)^1 (1-0.5)^{2} = 3 \cdot 0.5 \cdot 0.25 = \frac{3}{8}$
    \item $P(X=2) = \binom{3}{2} (0.5)^2 (1-0.5)^{1} = 3 \cdot 0.25 \cdot 0.5 = \frac{3}{8}$
    \item $P(X=3) = \binom{3}{3} (0.5)^3 (1-0.5)^{0} = 1 \cdot 0.125 \cdot 1 = \frac{1}{8}$
\end{itemize}

Como se puede observar, los resultados coinciden exactamente con los que obtuvimos en la Parte a) al contar los casos favorables sobre los casos posibles.

\vspace{0.5cm}

\section*{Anexo: ¿Importa el orden de nacimiento en este problema?}

Una duda conceptual muy frecuente en este tipo de ejercicios es si el orden en que nacen los hijos importa o no. La respuesta es que \textbf{depende de qué parte del problema estemos analizando}.

\begin{itemize}
    \item \textbf{Para definir el espacio muestral (casos posibles): SÍ importa el orden.} \\
    Los hijos nacen en una secuencia temporal (el mayor, el del medio, el menor). Si ignoráramos el orden por completo, podríamos cometer el error de pensar que solo hay 4 resultados posibles en nuestro universo: tener 0, 1, 2 o 3 varones, y creer erróneamente que todos son equiprobables (con probabilidad $1/4$ cada uno). Para que la matemática sea correcta, debemos contar cada secuencia exacta como un caso distinto. Así, tener ``Varón-Varón-Niña'' es un escenario distinto a tener ``Niña-Varón-Varón''. Esto nos revela la verdadera estructura: hay $2 \times 2 \times 2 = 8$ secuencias de nacimientos posibles y equiprobables.
    
    \item \textbf{Para la variable aleatoria $X$ (casos favorables): NO importa el orden.} \\
    Nuestra variable $X$ se define simplemente como la ``cantidad total de niños''. Funciona como un agrupador o contador final que ignora quién es el hermano mayor o el menor. Toma aquellas secuencias distintas que mencionamos antes (Varón-Varón-Niña, Varón-Niña-Varón, y Niña-Varón-Varón) y las mete todas juntas en la misma categoría: $X = 2$.
\end{itemize}

\end{document}
